\chapter{Introduction}

\textit{Convolutional neural networks}, often abbreviated as CNNs,
are inspired by the vertebrate visual cortex and play a central role in computer vision.
They are used to recognize objects, classify images, detect shapes, and automatically extract
visual information.

Their effectiveness stems largely from their ability to learn local patterns in
images and progressively combine them into more abstract representations.

\section*{Vulnerabilities and Adversarial Attacks}

Despite their strong performance, these models can exhibit structural vulnerabilities.
An image can be modified slightly, in a way that is imperceptible to the human eye,
yet causes a substantial change in the network's prediction. This phenomenon is known as an
\textbf{adversarial perturbation}~\cite{szegedy2014intriguing}.

A particularly interesting aspect of these perturbations is that they can sometimes be universal. In some cases,
a perturbation can be found that is based directly on the network's architecture
(how it operates), exploiting a structural weakness in its internal workings.
Once found, such a perturbation can be applied to any image,
\textbf{regardless} of its content, and cause the model to make incorrect
predictions on a substantial proportion of images. This is called a \textbf{universal} adversarial perturbation~\cite{moosavi2017universal}.

The central question is then where this sensitivity comes from. One possible explanation
attributes it to the training data or model optimization. In this work, we
explore another possibility: \textbf{the very structure of convolutional networks}. CNNs
rely on convolutional layers, which have specific mathematical properties.

\section*{The Fourier Basis Approach}

To investigate this question, we use the Fourier basis. An image can be
described pixel by pixel, but it can also be decomposed into frequency patterns: slow variations,
rapid oscillations, and horizontal, vertical, or diagonal patterns. These frequency patterns provide
another way to describe the structures present in an image.

Analyzing an image not only in the spatial domain, through its pixels, but
also in the frequency domain therefore helps us understand which image components influence
the network's behavior and which perturbation directions may become particularly sensitive.

This connection is the starting point of the work by Tsuzuku and Sato~\cite{tsuzuku2019structural}, who show
that certain convolutional networks are particularly sensitive to directions drawn from the Fourier basis.
Their idea is that some universal adversarial perturbations can be better understood by studying
how convolutional layers process frequencies.

\begin{center}
    \vspace{1em}
    \rule{0.4\textwidth}{0.4pt} \\
    \vspace{1em}
    \begin{minipage}{0.85\textwidth}
        \centering
        \large \itshape
        How does the convolutional architecture of CNNs give rise to preferred directions of vulnerability in
        the Fourier domain, and to what extent does this structural sensitivity help us understand
        the existence of universal adversarial perturbations?
    \end{minipage} \\
    \vspace{1em}
    \rule{0.4\textwidth}{0.4pt}
    \vspace{1em}
\end{center}

The goal of this report is to revisit this idea and reconstruct it step by step in greater detail,
introducing the concepts needed to understand the results so that any reader with a basic knowledge
of linear algebra and machine learning can understand the implications of this work.

\section*{Report Organization}

To address this question, the report is organized as follows:

\begin{itemize}[label=\textbullet, leftmargin=1.5cm]
    \item \textbf{Chapter 2:} General background on convolutional networks, robustness, and adversarial perturbations.
    \item \textbf{Chapter 3:} Theoretical foundations connecting convolution and the Fourier basis.
    In particular, this chapter shows how Fourier functions emerge naturally in linear analysis.
    \item \textbf{Chapter 4:} Experimental reproduction, the protocol used, and frequency sensitivity \textit{heatmaps},
    followed by a discussion of the results and limitations.
    \item \textbf{Conclusion:} Summary of the findings and directions for future research.
\end{itemize}
